% SLICEPARITYDEPTH1 2026-10-16 15:20 Asia/Seoul
\documentclass[11pt]{article}
\usepackage[T1]{fontenc}
\usepackage{lmodern,amsmath,amssymb,amsthm,booktabs,hyperref}
\usepackage[margin=0.9in]{geometry}
\newtheorem{theorem}{Theorem}
\newtheorem{lemma}[theorem]{Lemma}
\newtheorem{corollary}[theorem]{Corollary}
\newcommand{\EXACT}{\textbf{PROVED/EXACT}}
\newcommand{\KNOWN}{\textbf{KNOWN THEOREM}}
\newcommand{\COMP}{\textbf{COMPUTATIONAL}}
\newcommand{\OPEN}{\textbf{OPEN}}
\newcommand{\RET}{\textbf{FALSIFIED/RETIRED}}
\title{SLICEPARITYDEPTH1: A Constant-Weight Tseitin Embedding and the Depth--Size Frontier}
\author{P versus NP project}
\date{16 October 2026}
\begin{document}\maketitle

\noindent\textbf{Status.} The constant-weight Tseitin embedding, dual-rail restriction theorem, binary-depth lower bound, and cancellation-order counterexample are \EXACT at their declared models, pending independent review. H\aa stad's PARITY lower bound for constant-depth unbounded-fanin circuits is \KNOWN. The resulting bounded-depth monotone/set-circuit lower bounds are unconditional. No unrestricted superpolynomial monotone-circuit lower bound is proved; $P\ne NP$ remains \OPEN.

\section{IDEAS PROPOSED BEFORE LITERATURE}
\paragraph{A: parity restriction.} Reuse the fixed-support Tseitin signing family not as a full lower-bound witness, but as a restriction that forces any slice circuit to compute PARITY on the charge bits. Missing theorem: the restriction must have exact Hamming weight and literal/constant coordinate maps.
\paragraph{B: cancellation order as depth, not size.} Interpret the multilinear degree on the parity restriction as an exact interaction/cancellation order. Prediction: it yields a logarithmic binary-depth lower bound but not a superpolynomial size lower bound.
\paragraph{C: constant-depth lower-bound ladder.} Import H\aa stad's depth-$d$ PARITY lower bound to obtain an asymptotic size-depth curve for the exact slice.
\paragraph{D: use the result as a hostile control on the factorial-series program.} Falsifier for the naive series route: linear cancellation order may coexist with linear circuit size if depth is logarithmic.

\section{LITERATURE CHECK AND REVISIONS}
H\aa stad's switching-lemma lower bound gives, for every fixed $d\ge2$, size $\exp(\Omega_d(n^{1/(d-1)}))$ for depth-$d$ unbounded-fanin AND/OR/NOT circuits computing PARITY. Standard quantitative versions imply that polynomial-size PARITY circuits require depth $\Omega(\log n/\log\log n)$ in the unbounded-fanin model. The result applies a fortiori after input negations are pushed to literals.

The previous project report SLICEFINITERIGID1 proved that fixed-support Tseitin restrictions expose both literal polarities and therefore cannot transfer a known \emph{monotone} lower bound. That negative result is retained. The present observation is that exposing both polarities is exactly what is needed to invoke known \emph{general constant-depth} PARITY lower bounds.

SLICEASYMSERIES1 showed that the UNSAT probability is represented by cancellations extending to growing order. The present restriction demonstrates that growing interaction order alone is not a size lower bound: PARITY has degree $n$ but linear-size logarithmic-depth dual-rail circuits.

\section{CONSTRUCTION OR PROOF ATTEMPT}
\subsection{Exact-density Tseitin embedding}
Let $n\ge6$ be even and let $G_n=C_{n/2}\square K_2$ be the connected cubic prism graph. It has $n$ vertices and $3n/2$ edges. Associate one Boolean formula variable $x_e$ with each edge.

For a charge vector $b\in\{0,1\}^n$, the vertex $v$ imposes
\[\bigoplus_{e\ni v}x_e=b_v.\]
The parity equation on its three incident edge variables is encoded by the four proper 3-clauses forbidding the four wrong-parity local assignments. Let $T_b$ be the conjunction of all $4n$ vertex clauses.

\begin{lemma}[Tseitin parity]
$T_b$ is satisfiable iff $\bigoplus_v b_v=0$.
\end{lemma}
\begin{proof}
Summing all vertex equations over $\mathbb F_2$ cancels every edge variable twice, so odd total charge is inconsistent. Conversely, on a connected graph, delete one vertex equation; the remaining equations have full row rank $n-1$ over the incidence space, so every even charge vector is solvable.
\end{proof}

Set the declared target variable count to
\[q=3n,\qquad t=3q=9n.\]
The Tseitin core uses $3n/2$ target variables. Introduce $3n/2$ fresh variables and choose any fixed set $P_n$ of exactly $5n$ distinct proper 3-clauses on the fresh variables that is simultaneously satisfiable, for example clauses chosen from the pool satisfied by the all-one assignment. For all sufficiently large $n$ that pool has more than $5n$ members.

Define
\[F_b=T_b\wedge P_n.\]
Then every $F_b$ has exactly
\[4n+5n=9n=3q\]
distinct proper 3-clauses on the declared $q$ variables, and
\[\boxed{U_{q,3q}(F_b)=\operatorname{PARITY}(b).}\]

\subsection{Dual-rail coordinate theorem}
At a fixed vertex $v$, the two four-clause parity encodings partition the eight signed clauses on the three incident edge variables. Hence for every target clause coordinate $C$, its presence on the family $F_b$ is one of
\[0,\quad1,\quad b_v,\quad1-b_v.\]

\begin{theorem}[Literal restriction]
If a monotone AND/OR circuit $C$ computes $U_{q,3q}$, then substituting the above constants and literals into $C$ yields an ordinary De Morgan circuit of no larger size or depth computing PARITY on $n=q/3$ input bits. The same statement holds for a union/intersection set circuit computing the exact-layer UNSAT family, by evaluating membership of each intermediate family on $F_b$.
\end{theorem}
\begin{proof}
Every input star or selector becomes a literal or a constant on the family $\{F_b\}$. Union/intersection become OR/AND on the restricted membership bits. Correctness follows from the preceding parity identity.
\end{proof}

\subsection{Constant-depth size lower bounds}
\begin{corollary}[Depth-$d$ frontier]
For every fixed $d\ge2$, every depth-$d$ monotone circuit or depth-$d$ set circuit computing the exact-density UNSAT slice has size
\[\boxed{\exp\!\left(\Omega_d(q^{1/(d-1)})\right).}\]
In selector-input length $N=8\binom q3=\Theta(q^3)$ this is
\[\boxed{\exp\!\left(\Omega_d(N^{1/(3(d-1))})\right).}\]
\end{corollary}
\begin{proof}
Apply H\aa stad's PARITY lower bound to the restricted circuit on $n=q/3$ charge variables.
\end{proof}

Thus the lower-bound family forms an explicit asymptotic series over depth:
\[S_d(N)\ge\exp\!\left(\Omega_d(N^{1/(3(d-1))})\right),\qquad d=2,3,\ldots.\]
For standard quantitative versions of the switching-lemma theorem, a polynomial-size unbounded-fanin circuit on this slice must have depth at least
\[\Omega\!\left(\frac{\log q}{\log\log q}\right)=\Omega\!\left(\frac{\log N}{\log\log N}\right).\]

\subsection{Exact binary-depth lower bound from cancellation order}
Interpret a binary AND/OR circuit on the restricted family as a Boolean function of $b$. Replace each gate by its unique real multilinear polynomial on $\{0,1\}^n$. Input literals have degree one. For binary gates,
\[\deg(f\wedge g)\le\deg f+\deg g,\qquad\deg(f\vee g)\le\deg f+\deg g.\]
Hence a binary circuit of depth $h$ has degree at most $2^h$.

\begin{theorem}[Binary depth]
Every binary monotone/set circuit computing the slice has depth at least
\[\boxed{\lceil\log_2 n\rceil=\Omega(\log N).}\]
\end{theorem}
\begin{proof}
The unique multilinear polynomial of PARITY on $n$ bits has degree $n$, while the restricted circuit has degree at most $2^h$.
\end{proof}

\subsection{Cancellation-order route is insufficient for size}
Maintain a dual-rail pair $(P,\bar P)$ for the parity of a block. Combining two dual-rail parity pairs requires only a constant number of AND/OR gates:
\[P'=(P\wedge\bar Q)\vee(\bar P\wedge Q),\qquad
\bar P'=(P\wedge Q)\vee(\bar P\wedge\bar Q).\]
A balanced binary tree computes PARITY with $O(n)$ gates and $O(\log n)$ depth from the exposed literals $b_i,1-b_i$.

\begin{corollary}[Growing cancellation order is not a size lower bound]
Linear interaction degree, linear inclusion--exclusion order, or any analogous quantity that can double under a binary union/intersection level cannot by itself imply superpolynomial set-circuit size. On this exact slice restriction, the required order is $n$ while an $O(n)$-size circuit realizes it.
\end{corollary}

This directly retires the interpretation of SLICEASYMSERIES1's ``growing cancellation order'' as a standalone size lower bound. Such an order can still be useful only when accompanied by a no-mass-production theorem controlling how many independent high-order interactions one gate can create or share.

\section{TESTS AND OUTCOME}
A finite exhaustive audit used prism graphs with $n=6$ and $n=8$. The graphs had respectively 9 and 12 edge variables. For all $2^n$ charge vectors, exhaustive search over all edge assignments found zero mismatches with the theorem ``SAT iff total charge is even.'' At every vertex the two four-clause encodings were disjoint and their union contained all eight signed clauses on that support. The audit therefore checked 320 charge vectors and the complete local dual-rail partition with zero discrepancies. These checks are \COMP; the graph-parity and restriction statements above are symbolic.

\paragraph{Outcome.}
The exact slice now has an unconditional depth-size lower-bound ladder beyond depth two. This does not reach unrestricted circuit size, but it imposes a necessary structural regime on any hypothetical polynomial-size monotone/set circuit: it must exploit growing depth and DAG sharing. The infinite-series program is correspondingly narrowed: high cancellation order is necessary but not sufficient; the missing invariant must quantify mass production or independent interaction rank across many orders, not merely the maximum order attained.

\paragraph{Next target.}
Search for a gate-stable \emph{multiplicity} measure of high-order interactions that is large for the full UNSAT transversal family but remains small for the linear-size parity control. In parallel, attempt a selected-input Jukna construction whose respecting measures capture support expansion rather than raw cancellation order. Any successful invariant must survive Hiltgen--Paterson mass production and the dual-rail parity restriction above.

Root \texttt{main.tex} remains truncated and was not modified.

\begin{thebibliography}{4}
\bibitem{Hastad86} J. H\aa stad, \emph{Almost Optimal Lower Bounds for Small Depth Circuits}, STOC 1986.
\bibitem{Wegener85} I. Wegener, \emph{On the Complexity of Slice Functions}, TCS 38, 1985.
\bibitem{HP95} A. P. Hiltgen and M. S. Paterson, \emph{PI$_k$ Mass Production and an Optimal Circuit for the Ne\v{c}iporuk Slice}, Computational Complexity 5, 1995.
\bibitem{Tseitin68} G. S. Tseitin, \emph{On the Complexity of Derivation in Propositional Calculus}, 1968.
\end{thebibliography}
\end{document}
% FINAL_VERIFICATION_SENTINEL::MAIN::SLICEPARITYDEPTH1::2026-10-16_1520::P_VS_NP_OPEN